% !TEX program = lualatex
\documentclass[12pt]{ltjsarticle}
\usepackage{luatexja}
\usepackage{amsmath,amsthm,amssymb}
\usepackage{tikz}
\usetikzlibrary{arrows.meta, positioning}

\title{Q2: 微積分学の基本定理（右側導関数版）の数学的解説}
\author{TeXファイル作成: CODEX (OpenAI)\\Leanファイル生成: CODEX (OpenAI)}
\date{\today}

\theoremstyle{definition}
\newtheorem{definition}{定義}
\theoremstyle{plain}
\newtheorem{theorem}{定理}
\newtheorem{lemma}{補題}

\begin{document}
\maketitle

\section{問題設定}
本資料では、区間積分で定義される原始関数
\begin{equation}
F(t) = \int_{x = a}^{t} f(x)\,dx
\end{equation}
が連続関数 $f$ に対し $\mathrm{d}F/\mathrm{d}t = f(t)$ を満たすことを示す。これは Q2 に対応し、Lean では\newline
\texttt{integral\_hasDerivAt\_right}
を適用することで証明できる。

\section{数学的背景}
\subsection{仮定の整理}
\begin{itemize}
  \item $f\colon \mathbb{R} \to \mathbb{R}$ は連続。
  \item $a \in \mathbb{R}$ は固定。
  \item 区間 $[a, t]$ 上の積分はルベーグ積分を用いて定義する。
\end{itemize}
連続性から $f$ は任意の有界閉区間上で可積分であり、$F$ はよく定義される。

\subsection{定理の主張}
\begin{theorem}[微積分学の基本定理 (FTC-1)]
連続関数 $f$ に対して、$F(t) = \int_{a}^{t} f(x)\,dx$ とおくと、任意の $t \in \mathbb{R}$ で
\begin{equation}
\frac{dF}{dt}(t) = f(t)
\end{equation}
が成り立つ。
\end{theorem}

\subsection{証明のアイデア}
差商を用いて
\begin{equation}
\frac{F(t+h) - F(t)}{h} = \frac{1}{h}\int_{t}^{t+h} f(x)\,dx
\end{equation}
を評価する。連続性により $f$ が $t$ で一様連続に近い振る舞いを示すため、$h \to 0$ で右辺が $f(t)$ に収束する。
厳密には平均値の定理と $f$ の極限を用いて示す。

\section{Lean 定理との対応}
Lean の補題
\begin{center}
\texttt{intervalIntegral.integral\_hasDerivAt\_right}
\end{center}
は以下を仮定する：
\begin{enumerate}
  \item $f$ が $[a, t]$ 上で \texttt{IntervalIntegrable}。
  \item $f$ が $t$ で強可測かつ連続。
\end{enumerate}
連続性からこれらの前提が満たされ、結論として \texttt{HasDerivAt} を得る。Q2 の Lean 証明はこの定理を \texttt{simpa} で呼び出すだけで完了する。

\section{図式による全体像}
\begin{figure}[h]
  \centering
  \begin{tikzpicture}[node distance=2.8cm, every node/.style={rectangle, rounded corners, draw, align=center, font=\small}]
    \node (cont) {連続関数\\$f: \mathbb{R} \to \mathbb{R}$};
    \node (integrable) [right=of cont] {$f$ は $[a,t]$ 上で\\可積分};
    \node (primitive) [right=of integrable] {$F(t)=\int_{a}^{t} f$\\原始関数の定義};
    \node (lemma) [below=of integrable] {FTC-1 の差商評価\\(平均値の定理)};
    \node (lean) [below=of primitive] {Lean補題\\\texttt{integral\_hasDerivAt\_right}};
    \node (conclusion) [below=of lemma] {結論\\$\dfrac{dF}{dt}=f$};

    \draw[-{Latex[length=3mm]}] (cont) -- (integrable);
    \draw[-{Latex[length=3mm]}] (integrable) -- (primitive);
    \draw[-{Latex[length=3mm]}] (primitive) -- (lean);
    \draw[-{Latex[length=3mm]}] (lemma) -- (conclusion);
    \draw[-{Latex[length=3mm]}] (cont) -- (lemma);
    \draw[-{Latex[length=3mm]}] (lean) -- (conclusion);
  \end{tikzpicture}
  \caption{Q2 の論理構造}
\end{figure}

\section{まとめ}
Q2 の証明は、連続性が可積分性と強可測性を保証する点を把握し、FTC-1 をそのまま適用することで完結する。Lean では既存補題を活用することで、数学的議論を直接コード化できる。

\end{document}
